\begin{table}[t!]
    \centering
    \small
    \renewcommand\arraystretch{1.1}
    \begin{tabular}{  m{3.5cm} m{12cm}  }
        \hline
        \strut \textbf{Term} & \textbf{Description} \\
        \hline\hline
        Device & Any IoT device in the system. \par \textit{Example:} lights, locks, shades, TVs, thermostats, fridges. \\
        \hline
        Command & A message sent to, and associated action executed by, one IoT 
        device. \par \textit{Example:} close shades. \\
        \hline
        Routine & A sequence of commands. Triggered by time, or a set of triggering clauses, or manually. \par \textit{Example:} At 9pm, turn off all lights and lock all doors. \\
        \hline
        Smart Device ($S$) & A device that has (small) memory and computation capability. \par \textit{Examples:} fridges, thermostats, TVs. \\
        \hline
        Simple Device ($D$) & An IoT Device that is not a Smart Device, like a sensor or actuator. \par \textit{Examples:} lights, locks, shades. \\
        \hline \hline
    \end{tabular}
    \caption{\it \small Key Terms used in this paper, with examples for \Cref{comesh:fig:smart-building}.}
    \label{comesh:tab:terms}
\end{table}

\subsection{System Model}
\label{comesh:sec:system-model}

We assume a crash-recovery model: IoT devices may fail (crash) at any time, and then recover. We do not target Byzantine (malicious) faults. As is traditional in fault-tolerant literature, we assume that only up to $f$ simultaneous crash failures occur in the system. The value of $f$ can be calculated from deployment settings, e.g., in a smart building, by setting $f$ to the number of devices in the largest power domain. 

\Cref{comesh:tab:terms} summarizes key terms used in the paper. We assume a subset of devices ($\geq 3$) are smart devices, equipped 
with (small) memory and compute capability---other simple 
IoT devices may co-exist,
capable only of receiving commands, responding to queries, but not executing computations.

Since IoT devices are typically fixed early in the operation of a smart building or smart farm, 
we assume that the locations of devices are known (we perform mobility experiments later).
We assume an asynchronous system with unsynchronized clocks, where messages may be dropped.

Smart devices communicate via a mesh, i.e., a wireless ad-hoc network containing device-to-device Wifi links. Meshes are known to be robust even when large fractions of nodes fail~\cite{sengul2008adaptive}.
However, links can be lossy. An ad-hoc routing protocol \cite{olsrd}
exists among devices in this mesh.
We assume any smart device can send a command to any IoT device.  Finally, \sysname{} is layered atop a weakly-consistent membership service, e.g., Fuse~\cite{fuse}, or Medley \cite{medley}.  Such a membership protocol, running over an ad-hoc network,  maintains a (nearly) full membership list (of all alive IoT device IDs) at each \smart{}, delivering updates only eventually. Like Medley~\cite{medley}, \sysname{} assumes known device locations. Either precise (e.g., via a digital twin and map), or approximate device locations from localization algorithms work~\cite{han2013localization}.
